\section{Complexity of Formulae}

Whether a formula says the same thing inside a class as it does in $V$ depends on how its quantifiers range, so we need a way of measuring that. The first step is to give the quantifiers that range over a set, rather than over the whole universe, notation of their own.

% [CORRECTED] ``with $u$ and $v$ distinct variables'' is inserted, at the author's instruction --- without it, $\exists v \in v \, \varphi$ counts as bounded, but it stands for $\exists v \brac{v \in v \land \varphi}$, whose quantifier ranges over everything. Without Foundation that breaks $\Delta_0$-absoluteness in 3.4: $\exists v \in v \parenth{v = v}$ is false in the transitive class $\emptyset$ but true in $\set{a} \supseteq \emptyset$ if $a = \set{a}$, which $\ZFwoF$ does not rule out
\begin{boxnotation}[Bounded Quantifiers] \hfill
    \begin{itemize}
        \item $\forall u \in v \, \varphi$ stands for $\forall u \brac{u \in v \to \varphi}$

        \item $\exists u \in v \, \varphi$ stands for $\exists u \brac{u \in v \land \varphi}$
    \end{itemize}
    Quantifiers written this way, with $u$ and $v$ distinct variables, are said to be \textbf{bounded}.
\end{boxnotation}

A formula all of whose quantifiers are bounded can only ever look inside the sets it is handed, and such formulae sit at the bottom of the hierarchy we are about to build.

\begin{boxdefinition}[$\Delta_0$ Formulae]
    A formula is $\Delta_0$ iff every quantifier occurring in it is bounded.
\end{boxdefinition}

Everything above the bottom level is obtained by putting a block of unbounded quantifiers, all of the same kind, in front of a $\Delta_0$ formula.

\begin{boxdefinition}[$\Sigma_1$ and $\Pi_1$ Formulae] \hfill
    \begin{itemize}
        \item $\Sigma_1$ formulae look like
        \begin{align*}
            \exists u_0 \exists u_1 \cdots \exists u_n \, \varphi
        \end{align*}
        where $\varphi$ is $\Delta_0$.

        \item $\Pi_1$ formulae look like
        \begin{align*}
            \forall u_0 \forall u_1 \cdots \forall u_n \, \psi
        \end{align*}
        where $\psi$ is $\Delta_0$.
    \end{itemize}
\end{boxdefinition}

There are no $\Delta_1$ formulae. On the shapes above, a formula cannot begin with both an unbounded $\exists$ and an unbounded $\forall$; and allowing either block to be empty only puts the $\Delta_0$ formulae into both classes, which buys us nothing we did not already have. What there are is $\Delta_1^T$ formulae, and to get them we first have to fix a theory to do the translating for us.

\begin{boxdefinition}[Complexity Relative to a Theory]
    Let $T$ be a theory and let $\psi\of{\bar{u}}$ be a formula. We say $\psi\of{\bar{u}}$ is $\Delta_0^T$ iff there is a $\Delta_0$ formula $\varphi\of{\bar{u}}$ such that
    \begin{align*}
        T \vdash \forall \bar{u} \brac{\psi\of{\bar{u}} \leftrightarrow \varphi\of{\bar{u}}}
    \end{align*}
    Similarly for $\Sigma_1^T$ and $\Pi_1^T$. We say $\psi\of{\bar{u}}$ is $\Delta_1^T$ iff it is both $\Sigma_1^T$ and $\Pi_1^T$.
\end{boxdefinition}

Keep going to get $\Sigma_n$ and $\Pi_n$ formulae, and also $\Sigma_n^T$, $\Pi_n^T$ and $\Delta_n^T$ formulae.
